
\subsection{基础概念}
\begin{tabularx}{\columnwidth}{XX}
  \rowcolor{blue!30}
  \multicolumn{2}{l}{
    \parbox[c][2.5\ht\strutbox][c]{0.9\columnwidth}{\Large 基础概念}
  }\\
  体积（气体）&气体体积等于容器体积\\
  体积（液体）&液体体积由分子间距保持稳定\\
  压强定义&单位时间内对容器壁单位面积的平均冲击力\\
  温度定义&分子热运动的平均强度\\
\end{tabularx}

\textbf{微观解释}

体积、压强和温度都可以从分子热运动角度理解：气体分子无规则运动并不断碰撞容器壁，形成压强；温度反映分子热运动的剧烈程度。

\textbf{分子动理论}

物体由大量分子组成，分子在永不停息地做无规则运动，分子间同时存在引力和斥力。热现象的宏观规律由大量分子的统计行为决定。
